\documentclass[11pt,a4paper]{article}

% ---------------------------------------------------------
% PREAMBLE (stable, warning-free, Overleaf-safe)
% ---------------------------------------------------------
\usepackage[utf8]{inputenc}
\usepackage[T1]{fontenc}
\usepackage{lmodern}
\usepackage{amsmath, amssymb, amsfonts}
\usepackage{bm}
\usepackage{geometry}
\usepackage{graphicx}
\usepackage{hyperref}
\usepackage{cite}
\usepackage{authblk}
\usepackage{physics}
\usepackage{mathtools}

\geometry{margin=2.4cm}

\title{\textbf{Companion XLIV: Gravitational Slip in the Relaxation-Driven Cyclic Cosmology}}
\author{Michael Lehmann \\ \texttt{mi.lehmann@gmx.de}}
\date{April 2026}

\begin{document}
\maketitle

\begin{abstract}
The gravitational slip, defined as the difference between the two scalar metric 
potentials in the Newtonian gauge, provides a sensitive probe of the dynamical 
structure of spacetime. In the standard cosmological model, the slip is negligible 
in the absence of anisotropic stress, and the two potentials evolve identically. 
In the Relaxation-Driven Cyclic Cosmology (RDCC), the scale-dependent evolution of 
the gravitational potential introduces a natural gravitational slip, as the 
relaxation mechanism modifies the temporal evolution of the potentials in a 
scale-dependent manner. This Companion develops a continuous description of the 
gravitational slip in RDCC, deriving how the relaxation scale influences the 
metric potentials, the anisotropic stress, and the observational signatures in 
weak lensing, redshift-space distortions, and the integrated Sachs--Wolfe effect. 
The resulting framework provides a new geometric perspective on the RDCC Universe.
\end{abstract}
% ---------------------------------------------------------
\section{Introduction}

In the Newtonian gauge, scalar perturbations of the metric are described by two 
potentials, $\Phi$ and $\Psi$, which govern the motion of matter and the 
propagation of light. In the standard cosmological model, these potentials are 
equal in the absence of anisotropic stress, and their evolution is governed by the 
linear growth of structure.

In RDCC, the gravitational potential evolves differently. The relaxation mechanism 
suppresses the decay of small-scale modes and enhances the coherence of large-scale 
modes. This modifies the temporal evolution of the metric potentials in a 
scale-dependent manner, producing a natural gravitational slip. The slip becomes a 
direct tracer of the relaxation scale $k_{\mathrm{rel}}$ and the temporal structure 
of the RDCC gravitational field.
% ---------------------------------------------------------
\section{Metric Potentials in RDCC}

The perturbed metric in the Newtonian gauge is
\begin{equation}
    ds^{2}
    = -(1 + 2\Psi)\,dt^{2}
      + a^{2}(t)(1 - 2\Phi)\,d\mathbf{x}^{2}.
\end{equation}

In RDCC, the relaxation mechanism modifies the evolution of the potentials. The 
potential $\Phi$ evolves according to
\begin{equation}
    \Phi_{\mathrm{RDCC}}(k,a)
    = \Phi(k,a)
      \left[
        1 - \chi_{\Phi}
        \left(\frac{k}{k_{\mathrm{rel}}}\right)^{\alpha}
      \right],
\end{equation}
while the potential $\Psi$ evolves as
\begin{equation}
    \Psi_{\mathrm{RDCC}}(k,a)
    = \Psi(k,a)
      \left[
        1 - \chi_{\Psi}
        \left(\frac{k}{k_{\mathrm{rel}}}\right)^{\alpha}
      \right].
\end{equation}

Because the relaxation mechanism affects the two potentials differently, a 
gravitational slip arises naturally.
% ---------------------------------------------------------
\section{Definition and Evolution of the Gravitational Slip}

The gravitational slip is defined as
\begin{equation}
    \eta(k,a)
    = \frac{\Phi_{\mathrm{RDCC}}(k,a)}{\Psi_{\mathrm{RDCC}}(k,a)}.
\end{equation}

In the standard cosmological model, $\eta = 1$ in the absence of anisotropic 
stress. In RDCC, the slip becomes
\begin{equation}
    \eta_{\mathrm{RDCC}}(k,a)
    = \frac{
        1 - \chi_{\Phi}
        \left(\frac{k}{k_{\mathrm{rel}}}\right)^{\alpha}
    }{
        1 - \chi_{\Psi}
        \left(\frac{k}{k_{\mathrm{rel}}}\right)^{\alpha}
    }.
\end{equation}

This expression shows that the slip is scale-dependent and becomes most pronounced 
near the relaxation scale. The slip reflects the temporal structure of the RDCC 
gravitational field and provides a direct geometric signature of the relaxation 
mechanism.
% ---------------------------------------------------------
\section{Anisotropic Stress and the RDCC Geometry}

In general relativity, a gravitational slip can arise from anisotropic stress. In 
RDCC, the slip arises even in the absence of physical anisotropic stress, because 
the relaxation mechanism modifies the temporal evolution of the metric potentials 
in a scale-dependent manner.

The effective anisotropic stress can be defined as
\begin{equation}
    \pi_{\mathrm{eff}}(k,a)
    = \frac{k^{2}}{4\pi G a^{2}}
      \left[
        \Phi_{\mathrm{RDCC}}(k,a)
        - \Psi_{\mathrm{RDCC}}(k,a)
      \right].
\end{equation}

This effective stress encodes the geometric imprint of the relaxation mechanism and 
provides a new perspective on the structure of spacetime in RDCC.
% ---------------------------------------------------------
\section{Observational Signatures}

The gravitational slip affects several key observables. Weak lensing is sensitive 
to the combination $\Phi + \Psi$, while galaxy motions are sensitive to $\Psi$ 
alone. The slip therefore modifies the relation between lensing and clustering.

The integrated Sachs--Wolfe effect is sensitive to the time derivative of the 
metric potentials, and the slip introduces a scale-dependent modification to the 
ISW signal.

Redshift-space distortions depend on the velocity field, which is governed by 
$\Psi$. The slip therefore introduces a scale-dependent modification to the RSD 
signal.

These signatures provide a direct observational window into the relaxation scale 
and the geometric structure of the RDCC Universe.
% ---------------------------------------------------------
\section{Conclusion}

The gravitational slip in the Relaxation-Driven Cyclic Cosmology reflects the 
scale-dependent evolution of the metric potentials. The relaxation mechanism 
modifies the temporal evolution of $\Phi$ and $\Psi$ in a scale-dependent manner, 
producing a natural gravitational slip even in the absence of anisotropic stress. 
This slip affects weak lensing, redshift-space distortions, and the integrated 
Sachs--Wolfe effect, providing a direct geometric signature of the RDCC 
gravitational field. The present Companion establishes the theoretical foundation 
for future studies of the geometry of spacetime, the interplay between matter and 
light, and the observational tests of RDCC.

% ========================
% References
% ========================

\begin{thebibliography}{10}

\bibitem{Flagship} Lehmann, M. (2026). Relaxation-Driven Cyclic Cosmology (RDCC): A Minimal CPT-Symmetric Framework with a Single Parameter. Flagship Paper. Zenodo: \url{https://zenodo.org/records/19744958} (v43.0 or later).

\bibitem{Zenodo} The complete RDCC companion series (I--LVII) is available at the same Zenodo record above.

% Thematisch relevante Companions
\bibitem{CompXVI} Companion XVI: Boltzmann Code and Modified Hierarchy.
\bibitem{CompXXXI} Companion XXXI: RDCC Halo Model and Non-Linear Structure.
\bibitem{CompXXXII} Companion XXXII--XXXIX: Galaxy--Halo Connection, Cosmic Web, Lensing, RSD, ISW, tSZ, Rotation Curves, CGM.

% Wichtige externe Referenzen
\bibitem{Planck2020} Planck Collaboration (2020), Astron. Astrophys. 641, A6.
\bibitem{DESI2024} DESI Collaboration (2024/2025), arXiv: [aktuelle $f\sigma_8$/BAO-Paper].
\bibitem{JWST2024} Maiolino et al. (2024), Nature (JWST high-$z$ galaxies/quasars).
\bibitem{Kaiser1987} Kaiser, N. (1987), Mon. Not. R. Astron. Soc. 227, 1 (RSD).
\bibitem{Bartelmann2001} Bartelmann, M. \& Schneider, P. (2001), Phys. Rep. 340, 291 (Weak Lensing Review).
\bibitem{HuOkamoto2002} Hu, W. \& Okamoto, T. (2002), Astrophys. J. 574, 566 (CMB Lensing).
\bibitem{LewisChallinor2006} Lewis, A. \& Challinor, A. (2006), Phys. Rep. 429, 1 (CMB Lensing Review).
\bibitem{NFW1997} Navarro, J. F., Frenk, C. S. \& White, S. D. M. (1997), Astrophys. J. 490, 493 (Halo Profiles).

\end{thebibliography}

\end{document}
